\documentclass[12pt,a4paper]{article}
\usepackage[utf8]{inputenc}
\usepackage[T1]{fontenc}
\usepackage[spanish]{babel}
\usepackage{amsmath,amssymb}
\usepackage{geometry}
\usepackage{siunitx}
\usepackage{enumitem}
\usepackage{booktabs}
\usepackage{tcolorbox}
\usepackage{graphicx}
\geometry{margin=2.5cm}

\title{Simulacro --- Parcial 1\\Mecánica Aplicada}
\author{Mecánica Aplicada}
\date{11 de junio de 2026}

\newtcolorbox{solbox}[1][]{colback=green!5!white,colframe=green!50!black,fonttitle=\bfseries,title=Solución}

\begin{document}
\maketitle

\begin{center}
\textbf{Instrucciones:} Resuelva cada problema mostrando todo el desarrollo.
Duración estimada: 90 minutos.
\end{center}

\vspace{1cm}

% ============================================
\section*{Problema 1 (25\%) --- Coordenadas Cartesianas e Intrínsecas}
% ============================================

El movimiento de una partícula $P$ respecto a tierra está definido por:
\[
x(t) = 3t^2 - 6t \qquad\text{[m]},\qquad
y(t) = 4t^2 - 8t + 2 \qquad\text{[m]}.
\]

Para el instante $t = 1$ s, determine:

\begin{enumerate}[label=\alph*)]
    \item El vector velocidad y el vector aceleración.
    \item La componente tangencial y normal de la aceleración.
    \item El radio de curvatura de la trayectoria.
\end{enumerate}

\vspace{0.5cm}

\begin{solbox}
\textbf{Datos:} $x(t) = 3t^2 - 6t$, $y(t) = 4t^2 - 8t + 2$.

Derivamos:
\[
\dot{x} = 6t - 6, \quad \ddot{x} = 6
\qquad
\dot{y} = 8t - 8, \quad \ddot{y} = 8
\]

\textbf{a) Para $t = 1$ s:}

\[
\dot{x}(1) = 0,\quad \dot{y}(1) = 0 \Rightarrow \mathbf{v} = \mathbf{0}
\]
\[
\ddot{x} = 6,\quad \ddot{y} = 8 \Rightarrow \mathbf{a} = 6\mathbf{i} + 8\mathbf{j} \ \text{m/s}^2
\]

\textbf{b)} Como $\mathbf{v} = \mathbf{0}$, las componentes tangencial y normal no están
definidas en este instante (la partícula se detiene instantáneamente).
La aceleración $\mathbf{a}$ es puramente normal si la velocidad es cero
(cambio de dirección).

\textbf{c)} El radio de curvatura no está definido cuando $v=0$ pues $\rho = v^2/a_n$
da una indeterminación. Se requiere analizar el límite. Sin embargo,
podemos calcular $\rho$ a partir de:
\[
\rho(t) = \frac{|\mathbf{v}(t)|^3}{|\mathbf{v}(t) \times \mathbf{a}(t)|}
\]
Evaluando para $t$ cercano a $1$ s se obtendría un valor finito.
La trayectoria es una recta (ver inciso adicional), por lo que
$\rho \to \infty$ (trayectoria rectilínea).

\smallskip
\textbf{Verificación -- ecuación de trayectoria:}
Eliminamos $t$. De $x = 3t^2 - 6t$, $y = 4t^2 - 8t + 2$.

Observamos que $\frac{x}{3} = t^2 - 2t$ y $\frac{y-2}{4} = t^2 - 2t$.

Por lo tanto:
\[
\frac{x}{3} = \frac{y-2}{4} \quad\Rightarrow\quad y = \frac{4}{3}x + 2
\]
Es una recta. Para movimiento rectilíneo $\rho = \infty$.
\end{solbox}

\newpage

% ============================================
\section*{Problema 2 (25\%) --- Coordenadas Polares}
% ============================================

La partícula $P$ describe una trayectoria plana respecto a tierra. En un
instante dado se sabe que:
\[
r = 2\ \text{m},\quad \dot{r} = 3\ \text{m/s},\quad \ddot{r} = 1\ \text{m/s}^2,
\]
\[
\theta = 30^\circ,\quad \dot{\theta} = 4\ \text{rad/s},\quad \ddot{\theta} = -2\ \text{rad/s}^2.
\]

Determine:

\begin{enumerate}[label=\alph*)]
    \item El vector velocidad en componentes polares.
    \item El vector aceleración en componentes polares.
    \item La magnitud de la velocidad y la aceleración.
\end{enumerate}

\vspace{0.5cm}

\begin{solbox}
\textbf{Datos:} $r=2$, $\dot{r}=3$, $\ddot{r}=1$, $\dot{\theta}=4$, $\ddot{\theta}=-2$.

\textbf{a) Velocidad:}
\[
\mathbf{v} = \dot{r}\,\mathbf{e}_r + r\dot{\theta}\,\mathbf{e}_\theta
= 3\,\mathbf{e}_r + (2)(4)\,\mathbf{e}_\theta
= \boxed{3\,\mathbf{e}_r + 8\,\mathbf{e}_\theta \ \text{m/s}}
\]

\textbf{b) Aceleración:}
\[
\mathbf{a} = (\ddot{r} - r\dot{\theta}^2)\,\mathbf{e}_r
           + (2\dot{r}\dot{\theta} + r\ddot{\theta})\,\mathbf{e}_\theta
\]
\[
a_r = 1 - (2)(4^2) = 1 - 32 = -31\ \text{m/s}^2
\]
\[
a_\theta = 2(3)(4) + (2)(-2) = 24 - 4 = 20\ \text{m/s}^2
\]
\[
\boxed{\mathbf{a} = -31\,\mathbf{e}_r + 20\,\mathbf{e}_\theta \ \text{m/s}^2}
\]

\textbf{c) Magnitudes:}
\[
|\mathbf{v}| = \sqrt{3^2 + 8^2} = \sqrt{9 + 64} = \sqrt{73}
\approx \boxed{8.54\ \text{m/s}}
\]
\[
|\mathbf{a}| = \sqrt{(-31)^2 + 20^2} = \sqrt{961 + 400} = \sqrt{1361}
\approx \boxed{36.89\ \text{m/s}^2}
\]
\end{solbox}

\newpage

% ============================================
\section*{Problema 3 (25\%) --- Cinemática del Cuerpo Rígido}
% ============================================

La barra $AB$ de longitud $L = 2\ \text{m}$ está articulada a tierra en $A$.
El extremo $B$ se desliza sobre una superficie horizontal. En la configuración
mostrada, $AB$ forma $60^\circ$ con la horizontal, la velocidad angular de la
barra es $\omega_{AB} = 3\ \text{rad/s}$ en sentido antihorario y la
aceleración angular es $\alpha_{AB} = 2\ \text{rad/s}^2$ en sentido horario.

\begin{center}
\fbox{Nota: No hay figura disponible. Imagine la barra AB con A fijo, B sobre superficie horizontal.}
\end{center}

Determine:

\begin{enumerate}[label=\alph*)]
    \item El vector velocidad del punto $B$.
    \item El vector aceleración del punto $B$.
    \item La velocidad del punto medio $M$ de la barra.
\end{enumerate}

\vspace{0.5cm}

\begin{solbox}
\textbf{Datos:} $L = 2\ \text{m}$, $\omega = 3\ \text{rad/s}$ (antihorario, $+$),
$\alpha = -2\ \text{rad/s}^2$ (horario, $-$). $\theta = 60^\circ$.

Sistema de referencia: $\mathbf{i}$ horizontal derecha, $\mathbf{j}$ vertical arriba.

Vector de $A$ a $B$:
\[
\mathbf{r}_{AB} = L\cos\theta\,\mathbf{i} + L\sin\theta\,\mathbf{j}
= 2\cos60^\circ\,\mathbf{i} + 2\sin60^\circ\,\mathbf{j}
= \mathbf{i} + \sqrt{3}\,\mathbf{j} \ \text{m}
\]

\textbf{a) Velocidad de B:}
$A$ está fijo a tierra ($\mathbf{v}_A = \mathbf{0}$).
\[
\mathbf{v}_B = \mathbf{v}_A + \boldsymbol{\omega} \times \mathbf{r}_{AB}
= \mathbf{0} + (3\mathbf{k}) \times (\mathbf{i} + \sqrt{3}\mathbf{j})
\]
\[
\mathbf{v}_B = 3\mathbf{k} \times \mathbf{i} + 3\mathbf{k} \times \sqrt{3}\mathbf{j}
= 3\mathbf{j} - 3\sqrt{3}\,\mathbf{i}
= \boxed{-3\sqrt{3}\,\mathbf{i} + 3\mathbf{j} \ \text{m/s}}
\]

\textbf{b) Aceleración de B:}
\[
\mathbf{a}_B = \mathbf{a}_A + \boldsymbol{\alpha} \times \mathbf{r}_{AB}
             + \boldsymbol{\omega} \times (\boldsymbol{\omega} \times \mathbf{r}_{AB})
\]
$\mathbf{a}_A = \mathbf{0}$, $\boldsymbol{\alpha} = -2\mathbf{k}$.

Término angular:
\[
\boldsymbol{\alpha} \times \mathbf{r}_{AB} = (-2\mathbf{k}) \times (\mathbf{i} + \sqrt{3}\mathbf{j})
= -2\mathbf{k} \times \mathbf{i} -2\mathbf{k} \times \sqrt{3}\mathbf{j}
= -2\mathbf{j} + 2\sqrt{3}\,\mathbf{i}
\]

Término centrípeto:
\[
\boldsymbol{\omega} \times (\boldsymbol{\omega} \times \mathbf{r}_{AB})
= 3\mathbf{k} \times (3\mathbf{k} \times (\mathbf{i} + \sqrt{3}\mathbf{j}))
\]
Primero $\boldsymbol{\omega} \times \mathbf{r}_{AB} = \mathbf{v}_B$ (ya calculado).
\[
\boldsymbol{\omega} \times \mathbf{v}_B = 3\mathbf{k} \times (-3\sqrt{3}\,\mathbf{i} + 3\mathbf{j})
= -9\sqrt{3}\,\mathbf{k}\times\mathbf{i} + 9\,\mathbf{k}\times\mathbf{j}
= -9\sqrt{3}\,\mathbf{j} - 9\mathbf{i}
\]

Sumando:
\[
\mathbf{a}_B = (2\sqrt{3} - 9)\,\mathbf{i} + (-2 - 9\sqrt{3})\,\mathbf{j}
\]
\[
\boxed{\mathbf{a}_B \approx (-5.54\,\mathbf{i} - 17.59\,\mathbf{j})\ \text{m/s}^2}
\]

\textbf{c) Velocidad del punto medio M:}
\[
\mathbf{r}_{AM} = \frac{1}{2}\mathbf{r}_{AB} = 0.5\,\mathbf{i} + 0.5\sqrt{3}\,\mathbf{j}
\]
\[
\mathbf{v}_M = \boldsymbol{\omega} \times \mathbf{r}_{AM}
= 3\mathbf{k} \times (0.5\mathbf{i} + 0.5\sqrt{3}\mathbf{j})
= 1.5\mathbf{j} - 1.5\sqrt{3}\,\mathbf{i}
\]
\[
\boxed{\mathbf{v}_M = -1.5\sqrt{3}\,\mathbf{i} + 1.5\mathbf{j}\ \text{m/s}}
\]
\end{solbox}

\newpage

% ============================================
\section*{Problema 4 (25\%) --- Rodadura}
% ============================================

Un disco de radio $R = 0.5\ \text{m}$ rueda sin deslizar sobre una superficie
horizontal. El centro $C$ del disco se mueve con velocidad $v_C = 4\ \text{m/s}$
hacia la derecha y aceleración $a_C = 2\ \text{m/s}^2$ hacia la derecha.

Determine:

\begin{enumerate}[label=\alph*)]
    \item La velocidad angular $\boldsymbol{\omega}$ del disco.
    \item La aceleración angular $\boldsymbol{\alpha}$ del disco.
    \item El vector velocidad del punto $P$ en la periferia, ubicado en la
          posición más alta del disco.
    \item El vector aceleración del mismo punto $P$.
\end{enumerate}

\vspace{0.5cm}

\begin{solbox}
\textbf{Datos:} $R = 0.5\ \text{m}$, $v_C = 4\ \text{m/s}$ ($\rightarrow$),
$a_C = 2\ \text{m/s}^2$ ($\rightarrow$).

\textbf{a) Velocidad angular:}
Rodadura sin deslizamiento: $v_C = \omega R$.
\[
\omega = \frac{v_C}{R} = \frac{4}{0.5} = 8\ \text{rad/s}
\]
Sentido horario (para que el disco avance a la derecha):
\[
\boxed{\boldsymbol{\omega} = -8\mathbf{k}\ \text{rad/s}}
\]

\textbf{b) Aceleración angular:}
\[
\alpha = \frac{a_C}{R} = \frac{2}{0.5} = 4\ \text{rad/s}^2
\]
Mismo sentido que $\omega$ (acelera):
\[
\boxed{\boldsymbol{\alpha} = -4\mathbf{k}\ \text{rad/s}^2}
\]

\textbf{c) Velocidad del punto P (punto más alto):}
Vector de $C$ a $P$: $\mathbf{r}_{CP} = R\mathbf{j} = 0.5\mathbf{j}$.
\[
\mathbf{v}_P = \mathbf{v}_C + \boldsymbol{\omega} \times \mathbf{r}_{CP}
\]
\[
\mathbf{v}_P = 4\mathbf{i} + (-8\mathbf{k}) \times (0.5\mathbf{j})
= 4\mathbf{i} - 4(\mathbf{k}\times\mathbf{j})
= 4\mathbf{i} - 4(-\mathbf{i})
= 4\mathbf{i} + 4\mathbf{i}
\]
\[
\boxed{\mathbf{v}_P = 8\mathbf{i}\ \text{m/s}}
\]
(El punto más alto tiene el doble de la velocidad del centro.)

\textbf{d) Aceleración del punto P:}
\[
\mathbf{a}_P = \mathbf{a}_C + \boldsymbol{\alpha} \times \mathbf{r}_{CP}
             + \boldsymbol{\omega} \times (\boldsymbol{\omega} \times \mathbf{r}_{CP})
\]

Término angular:
\[
\boldsymbol{\alpha} \times \mathbf{r}_{CP} = (-4\mathbf{k}) \times (0.5\mathbf{j})
= -2(\mathbf{k}\times\mathbf{j}) = -2(-\mathbf{i}) = 2\mathbf{i}
\]

Término centrípeto:
\[
\boldsymbol{\omega} \times (\boldsymbol{\omega} \times \mathbf{r}_{CP})
= (-8\mathbf{k}) \times [(-8\mathbf{k}) \times (0.5\mathbf{j})]
\]
Calculamos $\boldsymbol{\omega} \times \mathbf{r}_{CP}$:
\[
(-8\mathbf{k}) \times (0.5\mathbf{j}) = -4(\mathbf{k}\times\mathbf{j}) = -4(-\mathbf{i}) = 4\mathbf{i}
\]
\[
(-8\mathbf{k}) \times (4\mathbf{i}) = -32(\mathbf{k}\times\mathbf{i}) = -32\mathbf{j}
\]

Sumando:
\[
\mathbf{a}_P = 2\mathbf{i} + 2\mathbf{i} - 32\mathbf{j}
= \boxed{4\mathbf{i} - 32\mathbf{j}\ \text{m/s}^2}
\]
\end{solbox}

\vfill

\begin{center}
\rule{0.5\textwidth}{0.5pt}
\small Fin del simulacro. Revise sus respuestas con los criterios del curso.
\end{center}

\end{document}
